\begin{table}[tbp]
\centering
\small
\caption{Application 3: policy 100 bp looser or tighter than the Taylor rule}
\label{tab:app3}
\begin{tabular}{llrrrr}
\toprule
& & RE & \multicolumn{3}{c}{Learning} \\
\cmidrule(lr){3-3}\cmidrule(lr){4-6}
Variable & Statistic & Looser & Looser & Tighter & Looser $+$ tighter \\
\midrule
Inflation & largest effect & +0.666 & +0.717 & $-$0.730 & $-$0.0138 \\
 & quarter & 2023Q3 & 2024Q1 & 2024Q1 & 2023Q4 \\
 & average effect & +0.366 & +0.413 & $-$0.415 & $-$0.0015 \\
\addlinespace
Output growth & largest effect & +1.322 & +1.382 & $-$1.380 & +0.0065 \\
 & quarter & 2023Q2 & 2023Q2 & 2023Q2 & 2024Q2 \\
 & average effect & +0.061 & +0.029 & $-$0.028 & +0.0005 \\
\addlinespace
Policy rate & largest effect & $-$1.735 & $-$1.817 & +1.815 & $-$0.0099 \\
 & quarter & 2023Q2 & 2023Q2 & 2023Q2 & 2024Q2 \\
 & average effect & $-$0.368 & $-$0.337 & +0.334 & $-$0.0025 \\
\addlinespace
Learners' share in force & largest effect & -- & +0.0244 & $-$0.0295 & $-$0.0053 \\
 & quarter & -- & 2024Q3 & 2024Q2 & 2024Q2 \\
\bottomrule
\end{tabular}
\par\smallskip
\begin{minipage}{\linewidth}\footnotesize
\textit{Notes:} Deviations from Application 1 in pp (annualised rates) when the policy rate is held 100 bp a year below (looser) or above (tighter) the Taylor rule in 2022Q3--2023Q2; forecast 2022Q3--2026Q2. Largest effect: the deviation with the largest absolute value, and its quarter; average effect: mean over the 16 quarters. Under RE the tighter run is the mirror image of the looser one (largest $|$looser $+$ tighter$|$ below $10^{-13}$). The policy rate includes the rule's response. Source: \texttt{app3\_policy.csv}.
\end{minipage}
\end{table}
